\begin{proposition}[C.2.2 (Proved): the untwisted intersection form]
Let $W_A$ be a sufficiently small compact tubular representative of the
smooth total space $\mathcal X^+$ of A.2.2 over a closed $u$-disc.
Its boundary is the mapping torus of Engel's \emph{untwisted} family
$X(u)$. It is not the six-sphere mapping torus $M_\Phi$.
Write $\mathcal L=H^2(S;\Z)$, $F=c_1\mathcal O_S(S_\infty)$,
$m=c_1(L)$, and $v=6P-O-3F$ in the conventions of A.4.1.
Let $H$ be the infinity-section class on $R$ and $A$ the exceptional
Cartier divisor pulled back to $R^+$. Then
\[
 H_4(W_A;\Z)=
 \frac{\mathcal L\oplus\Z H\oplus\Z A}{\Z(m-A)}.
\]
For $\alpha=(a,h,b)\in H^2(R^+;\Z)$ set
\[
 d(\alpha)=(I-T^*)a+h(m-F)+bF.
\]
The middle intersection form of $W_A$ is
\[
 B(\alpha,\beta)=-d(\alpha)\cdot d(\beta).
\]
Its rank is two, its signature is $+2$, and its two nonzero Smith
factors are $1,11$. On the image modulo its radical it is represented,
in the ordered basis $(v,m)$, by
\[
 \begin{pmatrix}6&1\\1&2\end{pmatrix}.
\]
\end{proposition}
\begin{proof}
The local semistable charts $xy=u$ give a tubular representative which
retracts to $Q^+$; its boundary is the nearby-fibre mapping torus. One
may equivalently choose a regular neighbourhood of $Q^+$ in the smooth
total space and round its boundary. The fibrewise plumbing identifies
that boundary with the inverse image of a small positive $u$-circle.
The normalization mapping cylinder and $H_3(S)=0$ identify its fourth
homology with the cokernel of
$H_4(S)\to H_4(R^+)$. In the divisor coordinates of A.4.2 the two
sections are $D'_0=H+m-A$ and $D_\infty=H$. Their difference is
$m-A$, proving the group presentation.

The normalization is an immersed complex hypersurface in $W_A$.
Its normal line has class $-D'_0-D_\infty$, because the whole central
divisor is the divisor of the base coordinate $u$. Put
$r_0\alpha=a+bF$ and $r_\infty\alpha=a-hm$; these are the
restrictions to the two sections. The self-intersection of two
four-cycles pushed forward from the normalization contains the normal
Euler term and the intersections between distinct local branches.
All transverse complex branch intersections have positive sign.
Consequently it is
\begin{align*}
 &-r_0\alpha\cdot r_0\beta-r_\infty\alpha\cdot r_\infty\beta\\
 &\quad+r_0\alpha\cdot T^*r_\infty\beta
       +T^*r_\infty\alpha\cdot r_0\beta\\
 &=-\bigl(r_0\alpha-T^*r_\infty\alpha\bigr)
       \cdot\bigl(r_0\beta-T^*r_\infty\beta\bigr).
\end{align*}
This is an intersection-of-cycles calculation in the actual smooth
fourfold, rather than an assignment of a form to an unconstructed
central fibre. The isometry property of $T$ and $T^*m=m-F$ give
the stated expression for $d$. Also $d(m-A)=0$, so it descends.

The translation formula in A.4.1 shows that the image of $d$ is
$\langle F,v,m\rangle$. The surjectivity of the three intersection
coordinates proved there shows equality integrally. Its radical is
$\Z F$, since $F^2=Fv=Fm=0$ and
$v^2=-6$, $m^2=-2$, $vm=-1$. The remaining form is the displayed
positive matrix, of determinant eleven and with an entry equal to
one. Integral row and column reduction gives Smith factors $1,11$.
These computations, including the full eleven-by-eleven form, are
independently checked in \path{code/e_push_C_attempts.py}.
\end{proof}

